
\section{Martingale Berry--Esseen Step}

We apply the martingale Berry--Esseen theorem to the scalar martingale
\(u^\top M_n^{\text{RR}}\).
% The Poisson decomposition writes
% \(W^{\text{RR}} = -n^{-1 / 2} M_n^{\text{RR}} + D_{2, n}^{\text{RR}}\), with
% \(D_{2, n}^{\text{RR}} = O(n^{-1 / 2})\). The bracket concentration gives
% \(|u^\top \langle M^{\text{RR}} \rangle_n u - n \, \sigma_n^{2, \text{RR}}(u)|\)
% in \(L_p\) at scale \(B(u) \sqrt{p \, n}\).

\textbf{Bounded increments.} For \(X_l \coloneqq u^\top \Delta M_l^{\text{RR}}\),
\begin{equation}
\begin{aligned}
|X_l| \le \varkappa(u),
\quad
\varkappa(u) \coloneqq 6 \, t_{\text{mix}} \, C_{\mathcal{Q}} \, \lVert  \varepsilon  \rVert_\infty \, \lVert  u  \rVert.
\end{aligned}
\label{eq:M-RR-incr}
\end{equation}
The conditional variances sum to
\begin{equation}
\begin{aligned}
V_n^2 \coloneqq \sum_{l = 2}^{n - 1} \sigma_l^2 = u^\top \langle M^{\text{RR}} \rangle_n u.
\end{aligned}
\end{equation}

\textbf{Variance lower bound.} Set \(s_n^2 \coloneqq n \, \sigma_n^{2, \text{RR}}(u)\). If
\begin{equation}
\begin{aligned}
n \, \alpha \, a \ge \frac{2 \, C_3 \, \lVert  u  \rVert^2}{\sigma^2(u)}
\end{aligned}
\label{eq:variance-lb-condition}
\end{equation}
then \(s_n^2 \ge n \sigma^2(u) / 2\).
% At \(\alpha = c \, n^{-1 / 2}\), \eqref{eq:variance-lb-condition} is satisfied
% for \(n \ge (2 C_3 \lVert  u  \rVert^2 / (c \, a \, \sigma^2(u)))^2\).
The trivial upper bound
\begin{equation}
\begin{aligned}
\sigma_n^{2, \text{RR}}(u) \le C_{\mathcal{Q}}^2 \lVert  \Sigma_{\varepsilon}^{(\text{M})}  \rVert \lVert  u  \rVert^2
\end{aligned}
\end{equation}
gives \(s_n^2 \le K^2(u) \, n\) with
\(K(u) \coloneqq C_{\mathcal{Q}} \, \lVert  \Sigma_{\varepsilon}^{(\text{M})}  \rVert^{1 / 2} \, \lVert  u  \rVert\).

We apply Lemma~\ref{lem:external-martingale-be} to this scalar martingale; the number of
increments is at most \(n\), and the increment bound is \(\varkappa(u)\).

\begin{theorem}\label{thm:M-RR-BE}
  \textbf{(Stationary martingale Berry--Esseen bound.)}
  Assume \textbf{UGE 1}, \(\pi(\varepsilon) = 0\), \(\lVert  \varepsilon  \rVert_\infty < \infty\),
  \(\sigma^2(u) > 0\), \(0 < \alpha\), and \(2 \alpha \le \alpha_\infty\). There exist
  constants \(C_{K, 1}(u), C_{K, 2}(u) > 0\) depending only on \(\lVert  u  \rVert\),
  \(\sigma(u)\), \(C_{\mathcal{Q}}\), \(t_{\text{mix}}\), \(\lVert  \varepsilon  \rVert_\infty\),
  \(\lVert  \Sigma_{\varepsilon}^{(\text{M})}  \rVert\), and the constants
  \(L_B(\varkappa(u)), C_1, C_2\) of Lemma~\ref{lem:external-martingale-be}, such that for every
  \(n \ge 3\) satisfying
  the variance lower-bound condition \eqref{eq:variance-lb-condition},
\begin{equation}
\begin{aligned}
  d_K \left(
    \frac{u^\top M_n^{\text{RR}}}{\sqrt{n} \, \sigma_n^{\text{RR}}(u)},
    \mathcal{N}(0, 1)
  \right)
    \le \frac{C_{K, 1}(u) \, \log^{3 / 4} n}{n^{1 / 4}}
     + \frac{C_{K, 2}(u) \, \log n}{\sqrt{n}}.
\end{aligned}
\label{eq:M-RR-BE}
\end{equation}
\end{theorem}

\textit{Proof.} Apply Lemma~\ref{lem:external-martingale-be} to
\((X_l)_{l = 2}^{n - 1}\) with partial sum \(u^\top M_n^{\text{RR}}\),
deterministic scale \(s_n^2 = n \, \sigma_n^{2, \text{RR}}(u)\), increment bound
\(\varkappa(u)\), and \(p = \left\lceil \log n \right\rceil\). This martingale array has at most \(n\)
increments, so the first term in Lemma~\ref{lem:external-martingale-be} is bounded using
\((2 n + 1) \log(2 n + 1)\).
% Use \(s_n^2 \in [n \, \sigma^2(u) / 2, \, K^2(u) \, n]\) from the
% variance lower and upper bounds.

\textbf{Term I.}
\begin{equation}
\begin{aligned}
L_B(\varkappa(u)) \, \frac{(2 n + 1) \log(2 n + 1)}{s_n^3}
  \le \frac{6 \sqrt{2} \, L_B(\varkappa(u))}{\sigma^3(u)} \, \frac{\log n}{\sqrt{n}}
  =: \frac{C^{\text{(I)}}(u) \, \log n}{\sqrt{n}}.
\end{aligned}
\label{eq:term-I}
\end{equation}

\textbf{Term III.} With \(a_p \coloneqq 2 p / (2 p + 1)\),
\begin{equation}
\begin{aligned}
s_n^{1 / (2 p + 1)} \le \max(1, K(u))^{1 / (2 p + 1)} \, n^{1 / (2 (2 p + 1))}.
\end{aligned}
\end{equation}
For \(p \ge \log n\),
\begin{equation}
\begin{aligned}
C_2 \, s_n^{- a_p} \, p \, \varkappa(u)^{a_p}
  \le \sqrt{2} \, e^{3 / 4} \, \frac{C_2 \, \max(1, \varkappa(u))}{\sigma(u)} \, \frac{p}{\sqrt{n}}
  \le \frac{C^{\text{(III)}}(u) \, \log n}{\sqrt{n}},
\end{aligned}
\label{eq:term-III}
\end{equation}
% The constant absorbs the bounded powers of \(K(u)\) and \(\varkappa(u)\).

\textbf{Term II.} By bracket concentration,
\begin{equation}
\begin{aligned}
(\mathbb{E} | V_n^2 - s_n^2 |^p)^{1 / p}
  \le B(u) \, \sqrt{p \, n},
\quad
B(u) \coloneqq C_4 \, C_{\mathcal{Q}}^2 \, \lVert  u  \rVert^2 \, \lVert  \varepsilon  \rVert_\infty^2 \, t_{\text{mix}}^{5 / 2},
\end{aligned}
\label{eq:Bu-def}
\end{equation}
The conditional-variance term is bounded by
\begin{equation}
\begin{aligned}
&C_1 \, \sqrt{p} \, s_n^{- a_p} \, (\mathbb{E} | V_n^2 - s_n^2 |^p)^{1 / (2 p + 1)} \\
&\quad\le C \, C_1 \, \sigma^{-1}(u)
        \, \max(1, K(u))^{1 / (2 p + 1)}
        \, \max(1, B(u))^{p / (2 p + 1)} \\
&\quad \quad \times p^{(3 p + 1) / (2 (2 p + 1))}
        n^{- p / (2 (2 p + 1))}.
\end{aligned}
\label{eq:term-II-1}
\end{equation}

For \(p = \left\lceil \log n \right\rceil\),
\(p^{(3 p + 1) / (2 (2 p + 1))} \le C \log^{3 / 4} n\) and
\(n^{- p / (2 (2 p + 1))} \le C n^{-1 / 4}\). Thus
% The \(K(u)\)- and \(B(u)\)-factors are absorbed into the direction-dependent
% constant.
\begin{equation}
\begin{aligned}
\text{Term II}
  \le \frac{C^{\text{(II)}}(u) \, \log^{3 / 4} n}{n^{1 / 4}}.
\end{aligned}
\label{eq:term-II}
\end{equation}

Adding the three bounds and setting
\(C_{K, 1}(u) \coloneqq C^{\text{(II)}}(u)\), \(C_{K, 2}(u) \coloneqq C^{\text{(I)}}(u) + C^{\text{(III)}}(u)\)
proves the martingale Berry--Esseen bound. \(\square\)

\begin{corollary}\label{cor:M-RR-BE-sigma}
  \textbf{(Stationary martingale asymptotic-normalization bound.)}
  Under the hypotheses of the previous theorem,
\begin{equation}
\begin{aligned}
  d_K \left(
    \frac{u^\top M_n^{\text{RR}}}{\sqrt{n} \, \sigma(u)},
    \mathcal{N}(0, 1)
  \right)
    \le \frac{C_{K, 1}(u) \, \log^{3 / 4} n}{n^{1 / 4}}
     + \frac{C_{K, 2}(u) \, \log n}{\sqrt{n}}
     + \frac{C_3 \, \lVert  u  \rVert^2}{n \, \alpha \, a \, \sigma^2(u)}.
\end{aligned}
\label{eq:M-RR-BE-sigma}
\end{equation}
% At \(\alpha = c \, n^{-1 / 2}\) the last term is \(O(n^{-1 / 2})\), hence
% absorbed into \(C_{K, 2}(u) \, \log n / \sqrt{n}\) up to a constant.
\end{corollary}

\textit{Proof.} Set \(r \coloneqq \sigma_n^{\text{RR}}(u) / \sigma(u)\) and
\(W \coloneqq u^\top M_n^{\text{RR}} / (\sqrt{n} \sigma_n^{\text{RR}}(u))\). Then the statistic
normalised by \(\sigma(u)\) is \(W r\). Under the variance lower-bound condition \eqref{eq:variance-lb-condition},
\(r \ge 1 / \sqrt{2}\), while the trivial upper bound on
\(\sigma_n^{2, \text{RR}}(u)\) gives \(r \le r_{\max}(u) < \infty\). On this compact
interval the standard normal cdf satisfies
\begin{equation}
\begin{aligned}
\sup_x |\Phi(x / r) - \Phi(x)| \le C_\Phi \, |r - 1|,
\quad
C_\Phi \coloneqq \sqrt{2} / \sqrt{\pi e}.
\end{aligned}
\end{equation}
Therefore
\begin{equation}
\begin{aligned}
d_K \left(W r, \mathcal{N}(0, 1)\right)
  \le d_K \left(W, \mathcal{N}(0, 1)\right) + C_\Phi \, |r - 1|.
\end{aligned}
\end{equation}
The variance comparison of Section 4.5 gives
\begin{equation}
\begin{aligned}
|r - 1|
  = \frac{|\sigma_n^{2, \text{RR}}(u) - \sigma^2(u)|}{\sigma(u) \, (\sigma_n^{\text{RR}}(u) + \sigma(u))}
  \le \frac{C_3 \, \lVert  u  \rVert^2}{n \, \alpha \, a \, \sigma^2(u)}.
\end{aligned}
\end{equation}
Adding this perturbation to the martingale Berry--Esseen bound proves the
stated claim after absorbing \(C_\Phi\) into the constants. \(\square\)

% This theorem concerns only the martingale term; the stationary assembly below
% adds the Poisson remainder and the Levin depth-two misadjustment.
